% ===========================================================================
%  08 — 结论
%
%  事实来源：paper_context/paper_outline.md §2.8
%  证据：     paper_context/claims_registry.md
%  术语：     paper_context/terminology.md
% ===========================================================================
\chapter{结论}\label{ch:conclusion}

本论文旨在设计、实现并从技术上验证一个基于遥测数据基础的教学运行时系统，用于高保真飞行仿真器中的程序性学习。所交付的贡献包括：一个可运行的双平台运行时、一个具有可推广设计原则的结构化视觉事实本体、一个附带基准测试证据的微调 VLM 适配器、一条可复现的数据流水线、跨骨干泛化结果，以及一个可扩展的评估基础设施，并附有通往人因研究的记录路径。本章总结已完成的工作，重申第\ref{ch:introduction}章中修订层次结构下的五项贡献，并概述未来工作方向。

% ===========================================================================
\section{已完成工作总结}\label{sec:conc-summary}

本论文产生了五项贡献，按创新性排序，并均有已验证声明和代码仓库工件的支撑。

\textbf{第一，结构化视觉事实本体设计，附抽象层次设计原则。}
13-fact v2 本体将座舱显示页面状态分解为跨三个显示区域（Left DDI、Right DDI、AMPCD）的离散、可局部验证的命题。该本体从 8-fact v1 基线演进而来，后者混合了抽象层次（例如复合目标 \texttt{ins\_go}），在完成类事实上产生了高达每 100 个留出样本 70 个的系统性假阳性。v2 分解为原子级视觉证据——将 \texttt{ins\_go} 拆分为 \texttt{ins\_grnd\_alignment\_text\_visible} 和 \texttt{ins\_ok\_text\_visible}，并将 \texttt{fcsmc\_complete} 分解为三个顺序子状态——在 Run-002 留出集上将关键假阳性减少了 64\%（13$\rightarrow$4），在 Run-004 留出集上减少了 89\%（70$\rightarrow$8）。本论文将该发现归纳为：VLM 应提取可观察的视觉原子；程序性解释应归于下游推理。
% 证据：C-VLM-01, C-VLM-02, C-VLM-03, C-VLM-04, C-VLM-05

\textbf{第二，完整的小数据 VLM 领域适配流水线。}
该流水线构成闭环：DCS 截图捕获、AI 预标注、Label Studio 中的人工审查、双语（英文+中文）监督微调导出、通过 Unsloth + PEFT + TRL 进行 LoRA 微调，以及在独立留出集上进行自动化基础模型与 LoRA 对比基准测试，并验证与训练集零重叠。该流水线以 180–220 个人工审查样本演示，产出 928 条双语训练数据。流水线中的每个工件均存储在代码仓库中，可独立检查和复现。
% 证据：C-VLM-03a, C-VLM-07, C-DATA-01, C-TRAIN-01

\textbf{第三，集成遥测-视觉-知识架构，具备可验证证据溯源。}
教学运行时强制要求每个叠加目标和步骤推荐都携带可追溯的证据引用（遥测变量、门控规则、RAG 片段或视觉事实）。证据约束通过运行时生成的 JSON Schema 在模式层面执行，该 Schema 包含 \texttt{if/then} 规则，在每个可操作输出被派发至仿真器之前对其进行验证——而不仅仅是为事后审计做日志记录。运行时编排完整的帮助循环：它从规范化遥测、增量摘要、候选步骤、确定性提示、叠加目标允许列表、检索到的知识片段以及视觉事实观察中组装结构化上下文；将该上下文传递给 LLM 以生成结构化帮助响应；依据模式、允许列表和证据溯源约束验证响应；执行经过验证的叠加操作；并将每个阶段记录为带有版本和时间戳的事件，以便离线回放。确定性降级路径确保当没有模型可用或模型输出验证失败时系统仍能正常运行。所有程序定义、门控规则、遥测映射和视觉事实本体均以声明式 YAML 配置（程序包）的形式表达，使系统可扩展到新的程序和仿真器。
% 证据：C-ARCH-01, C-ARCH-02, C-ARCH-03, C-PROC-01, C-PROC-02, C-TELEM-01, C-RAG-01, C-RUNTIME-01

\textbf{第四，跨骨干泛化证据。}
相同的双语 SFT 方案从 Qwen3.5-9B 迁移到 Gemma-4-31B 并获得一致改进，证实数据方案和本体设计并非特定于某一架构。Gemma LoRA 适配器分别实现了 0.9492 和 0.9715 的事实准确率（基础模型：0.8169 和 0.8146），在 Run-002 留出集上实现零关键假阳性，揭示了骨干层面在错误风格（保守 vs. 自信）上的差异，这对安全关键型部署决策具有重要参考价值。
% 证据：C-VLM-06, C-VLM-03b

\textbf{第五，可扩展的运行时和评估基础设施。}
本论文交付了一个程序包驱动、仿真器无关的架构，包含程序状态机、门控规则引擎、确定性降级、JSONL 事件日志、回放验证、自动化评分，以及一个用于基于 VR 的人因评估的设计层面框架——全部通过自动化测试实现。运行时同时支持桌面/DCS 和 VR（Quest~3）配置。已规划的人因评估设计——一项三条件受试者间研究（视频加笔记、无状态纯文本 LLM、基于状态证据教学），涵盖性能、工作负荷、学习效果、保持和迁移等多维度指标——记录于第\ref{sec:eval-planned}节。评估设计及执行所需的配套基础设施升级，构成了研究计划的下一阶段。
% 证据：C-ARCH-01, C-PROC-01, C-RUNTIME-01, C-HUMAN-01, C-HUMAN-02, C-HUMAN-03, C-DATA-COLLECT-01

\paragraph{本论文证明了什么。}
本工作确立了：一个具备结构化视觉事实提取功能的、基于遥测数据与证据约束的教学运行时在技术上是可行的，并且可以通过系统性、可复现的自动化基准测试进行验证。它证明了适量的人工审查数据（180–220 个样本）足以将通用 VLM 适配到特定领域的结构化提取任务，并在独立留出集上达到接近上限的准确率。它表明视觉事实本体的设计——具体而言，将高层次程序状态分解为低层次、可局部验证的原子这一决策——对模型可靠性具有可度量、可量化的影响。它提供了证据表明这些改进可跨 VLM 骨干泛化，证实驱动增益的是数据方案和本体设计，而非任何单一模型架构。并且，它交付了一个可运行的运行时基础设施，包含自动化测试、回放能力，以及通向受控人因评估的已记录扩展路径。

\paragraph{范围与边界。}
本论文未报告人类参与者的学习效果数据，也未声称该系统能够改善程序性技能获取、降低工作负荷或加速培训（相对于其他教学方法）。在本论文范围内，验证通过自动化基准测试进行；人因评估在已规划研究设计（第\ref{sec:eval-planned}节）中处理，属于未来工作。VR 运行时代码路径已实现（VR\_MIRROR、教学文本显示、复合面板捕获），但其与完整教学循环的集成尚未在开发环境之外进行演练。这些边界条件不限制 VLM 适配流水线所报告基准测试结果的内部有效性。

% ===========================================================================
\section{未来工作}\label{sec:conc-future}

本论文的贡献为若干未来工作方向开辟了道路，每个方向都直接建立在已完成的贡献之上。

\paragraph{VR 实验验证。}
最直接的扩展是使用人因受试者对 VR 运行时代码路径进行实验验证。VR 代码路径已实现（VR\_MIRROR 配置、教学文本 UDP 注入、VR 模式下的复合面板捕获）；剩下需要做的是在 DCS-VR 环境中通过研究参与者验证遥测流水线和帮助循环，并集成眼动和手部追踪输入以触发帮助。现有的端口-适配器架构支持这一点：核心教学逻辑是平台无关的，VR 输入/输出适配器可以在不修改核心的情况下进行验证。

\paragraph{人因评估。}
第\ref{sec:eval-planned}节描述的已规划评估设计应作为受控受试者间研究加以执行。这需要完成第\ref{sec:eval-planned}节确定的运行时数据采集升级（参与者/会话模式、NASA-TLX 集成、测验关联、匿名化流水线），从初级到中级人群中招募参与者，并执行完整的前测/培训/后测/保持协议。只有这项研究才能确定基于状态证据的教学方法是否在程序性学习效果、工作负荷和保持方面相对于基线条件产生可测量的改进。

\paragraph{本体向其他程序的扩展。}
13-fact v2 本体特定于 F/A-18C 冷启动任务及其三显示器复合面板布局。将本体扩展到 DCS 内的其他程序——如导航、武器使用或应急检查单——将检验视觉事实分解策略的通用性，并需要新的数据集采集和标注周期。程序包格式专为容纳此类扩展而设计：一个新程序包目录，包含其自身的 \texttt{vision\_facts.yaml}、\texttt{vision\_layout.yaml}、\texttt{step\_registry.yaml} 和 \texttt{telemetry\_map.yaml}，将定义一个新的训练场景而无需更改流水线工具。

\paragraph{多帧时序推理。}
当前的 VLM 流水线在单帧对上操作（预触发帧和触发帧），并在每个观察窗口上进行独立预测。在多个帧之间融入时序上下文可以提高涉及瞬态或状态转换的事实的准确率，例如 FCS-MC 测试序列（\texttt{fcsmc\_in\_test\_visible} $\rightarrow$ \texttt{fcsmc\_final\_go\_result\_visible}）。可能的方法包括将帧堆叠作为 VLM 的多图像输入、在滑动窗口上对事实预测进行递归聚合，或者使用一个轻量级时序模型，消费逐帧事实向量并生成时序一致的序列。

\paragraph{在线适配。}
当前的 LoRA 适配器在固定数据集上离线训练并作为静态模型部署。在线适配循环——系统在实时教学会话中采集新的标注帧（经导师或专家验证）并定期更新适配器——可以逐步提高边缘情况的准确率，并适应 DCS 版本或飞机变体间座舱显示配置的变化。这将需要能够与教学运行时并发运行的轻量级微调基础设施，以及对新采集标注的质量控制机制。

\paragraph{确定性程序引擎增强。}
当前的门控规则引擎使用一个包含四个操作符（\texttt{var\_gte}、\texttt{arg\_in\_range}、\texttt{flag\_true}、\texttt{time\_since}）的 v1 DSL。扩展 DSL 以添加更多操作符（如逻辑组合、序列计数器、多变量条件）并支持用户自定义门控函数，将提高程序包的表达能力，并减少对边缘情况手动步骤分类的依赖。

本论文所完成的工作确立了：一个具备结构化视觉事实提取功能的、基于遥测数据与证据约束的教学运行时在技术上是可行的，并可通过自动化基准测试进行验证。下一阶段的研究必须确定这一技术能力能否转化为沉浸式仿真环境中人类学习的实际改善。
